\documentclass[11pt]{article}
\usepackage[margin=1in]{geometry}
\usepackage{amsmath,amssymb,amsthm,mathtools}
\usepackage[colorlinks=true,linkcolor=blue,citecolor=blue,urlcolor=blue]{hyperref}
\newtheorem{theorem}{Theorem}
\newtheorem{lemma}[theorem]{Lemma}
\newcommand{\LL}{\mathcal L}
\newcommand{\Id}{\operatorname{Id}}
\newcommand{\tr}{\operatorname{tr}}
\newcommand{\str}{\operatorname{str}}
\newcommand{\Hol}{\operatorname{Hol}}
\title{Counterexample to the Fried conjecture}
\author{}
\date{11 October 2026\\[2pt]\small Manuscript draft}
\begin{document}
\maketitle

\begin{abstract}
We construct negatively curved metrics on a closed three-manifold and
acyclic unitary line characters for which the value at zero of the Ruelle
zeta function, divided by the squared analytic torsion, tends to zero.
The construction combines the hyperbolic twisted resonance calculation
with a conformal variation that moves a one-form resonance through zero.
Uniform holomorphy in the character parameter and a finite-cut torsion
variation formula determine the zeta value at the crossing. The argument
concerns the ordinary meromorphic zeta function and does not identify its
specialized value with a continued value from nearby metrics.
\end{abstract}

\section{Statement}
For a closed oriented negatively curved manifold $(\Sigma,g)$ and a
unitary representation $\chi$, put
\begin{equation}\label{zeta}
 Z_{g,\chi}(s)=\prod_{\gamma\ \mathrm{primitive}}
       \det\bigl(\Id-\chi(\gamma)e^{-s\ell_g(\gamma)}\bigr).
\end{equation}
The product is over primitive oriented geodesic-flow trajectories and
extends meromorphically. The Fried conjecture, in the form of
\cite[Conjecture 9]{Zw}, asserts that for acyclic $\chi$ and
$n=\dim\Sigma>2$,
\begin{equation}\label{fried}
 |Z_{g,\chi}(0)|^{(-1)^{n-1}}=T_\chi(\Sigma)^2,
\end{equation}
where $T_\chi$ is Ray--Singer analytic torsion. The survey's coordinate
is $\zeta_\chi(\lambda)=Z_{g,\chi}(-i\lambda)$.

\begin{theorem}\label{main}
There exist a closed oriented hyperbolic three-manifold $(\Sigma,g_0)$,
smooth negatively curved metrics $g_j\to g_0$, and acyclic unitary line
characters $\chi_j\to\mathbf1$ such that $Z_{g_j,\chi_j}$ is holomorphic
at zero and
\begin{equation}\label{ratio}
 \frac{|Z_{g_j,\chi_j}(0)|}{T_{\chi_j}(\Sigma)^2}\longrightarrow0.
\end{equation}
Consequently \eqref{fried} fails for all sufficiently large $j$.
\end{theorem}

The characters are acyclic at every member of the sequence; only their
limit is nonacyclic. The mechanism is a zero-resonance crossing at which
spectral specialization and continuation of nearby zero values need not
agree. We use the resonance calculations of Dang--Guillarmou--Rivi\`ere--Shen
(DGRS) \cite{DGRS}, the conformal variation of
Ceki\'c--Delarue--Dyatlov--Paternain (CDDP) \cite{CDDP}, and the dynamical
torsion of Chaubet--Dang \cite{CD}.
The resonance-crossing strategy was proposed in the conditional formal
reduction \cite{Crossing}. We derive the geometric and analytic input
below, keeping a fixed spectral cut through the crossing.

\section{Geometric and spectral preparation}
Choose $(\Sigma,g_0)$ with $b_1(\Sigma)=1$. Such a manifold can be
constructed as follows. A genus-two mapping class with symplectic
homology action
\[
 T=\begin{pmatrix}A&-\Id_2\\ \Id_2&0\end{pmatrix},\qquad
 A=\begin{pmatrix}2&1\\1&-1\end{pmatrix},
\]
has characteristic polynomial $p(x)=x^4-x^3-x^2-x+1$. Modulo two,
$p=\Phi_5$ is irreducible; a real root between one and two excludes
cyclotomicity, and the odd-degree coefficients exclude $p=q(x^k)$,
$k>1$. The homological pseudo-Anosov criterion
\cite[Propositions 1--2]{MS} and fibered hyperbolization \cite{Thurston}
therefore give a closed hyperbolic mapping torus. The Wang sequence and
$p(1)=-1$ give $b_1=1$.

Let $\omega\ne0$ be a real harmonic one-form and set
\begin{equation}\label{family}
 d_z=d+iz\omega\wedge,\qquad g_\tau=e^{-2\tau b}g_0.
\end{equation}
For real $z=\theta$, this connection defines a unitary character
$\chi_\theta$. All sufficiently small nonzero $\theta$ are acyclic.
Indeed, the small scalar and one-form Hodge spectral projections each
have rank one. Nontrivial holonomy makes the scalar eigenvalue
$\mu_0(\theta)$ positive, and $d_\theta f_\theta$ occupies the unique
small one-form slot at that same eigenvalue. Unitary Hodge duality
excludes harmonic forms in degrees two and three. This also leaves a
uniform positive gap in the coexact one-form spectrum. Since
$d^*\omega=0$, the first scalar eigenfunction correction vanishes, so
\begin{equation}\label{c0}
 \mu_0(\theta)=\kappa\theta^2+O(\theta^4),\qquad
 c_0(\theta):=-1+\sqrt{1-\mu_0(\theta)}
       =-\frac\kappa2\theta^2+O(\theta^4),\qquad
 \kappa=\frac{\|\omega\|_2^2}{\operatorname{vol}(\Sigma)}>0.
\end{equation}

Identify the unit bundles with $M=S_{g_0}\Sigma$ by radial scaling.
Write $\alpha_\tau,X_\tau$ for the contact form and Reeb field, and
\[
 L_{z,\tau}=\LL^{d_z}_{X_\tau},\qquad B_{z,\tau}=-L_{z,\tau}.
\]
Thus the spectral coordinate is $s$, with resolvent $(L+s)^{-1}$.
Fix a small circle $\gamma$ about zero containing no boundary resonance
at $(z,\tau)=(0,0)$ in any degree. Denote its horizontal generalized
resonance spaces by $U^k$ and its full spaces by $C^k$. Contact splitting
and Lefschetz give
\begin{equation}\label{ranks}
 \dim U^\bullet=(1,2,4,2,1),\qquad
 \dim C^\bullet=(1,3,6,6,3,1),\qquad
 C^k=U^k\oplus\alpha\wedge U^{k-1}.
\end{equation}

\begin{lemma}\label{axes}
On the hyperbolic axis, the horizontal spectra in degrees zero, one,
and two are respectively
\begin{equation}\label{hyper}
 (c_0),\qquad(c_0,c_0),\qquad(c_0,c_0,c_0,-c_0).
\end{equation}
There is a smooth real $b$ for which the untwisted one-form companion
root satisfies $t(0,\tau)=r\tau+o(\tau)$ with $r>0$. For small
$\tau>0$, the untwisted horizontal zero multiplicities are $(1,1,3,1,1)$.
\end{lemma}
\begin{proof}
Use DGRS's Selberg parameter $s$, centered at one. Bunke--Olbrich's
Selberg determinant formula \cite[Proposition 4.7]{BO} gives, up to
nonvanishing factors,
\[
 Z^D_{\mathbf1}(s)\doteq
 \det\bigl(\Delta_\theta^{(0)}-1+(s-1)^2\bigr),\qquad
 Z^D_{\nu_1}(1+w)\doteq
 \frac{\det(\Delta_\theta^{(1)}+w^2)}
      {\det(\Delta_\theta^{(0)}+w^2)},
\]
where $\nu_1$ is the pair of $SO(2)$ weights $\pm1$. This pair
contributes one graded determinant. The exact one-form spectrum cancels
the scalar denominator, including $\mu_0$; the remaining coexact spectrum
has the gap proved above. At $\theta=0$ the harmonic factors cancel
because $b_0=b_1=1$.

In DGRS (7.6), the shifts are $2(j-l)+p+2+k$, with $\sigma_p$ the
\emph{trace-free} symmetric representation. The factors with shifts at
most two reduce, near zero, to
\[
 P_0\doteq Z^D_{\mathbf1}(s+2),\qquad
 P_1\doteq Z^D_{\mathbf1}(s+2)^2,\qquad
 P_2\doteq Z^D_{\mathbf1}(s)Z^D_{\mathbf1}(s+2)^3.
\]
The other weight-two factor has virtual lift
$\operatorname{Sym}^2_0-\Lambda^1$ and shifted operators
$\nabla^*\nabla-2+2s+s^2$ and $\nabla^*\nabla+2+2s+s^2$.
DGRS Proposition 7.7's Bochner bound $\nabla^*\nabla\ge3$ on
$\operatorname{Sym}^2_0$ excludes its small roots; the unitary twist adds
no curvature term. Higher shifts lie in the zero-free Euler half-plane.
The scalar equations $\mu_0+s^2\pm2s=0$ now give \eqref{hyper}, with
algebraic multiplicities as in DGRS. We use BO only for its Selberg
formula, not its Ruelle normalization.

For the metric direction, normalize $u\in\mathcal C_\psi$ by
$\pi_*(d\alpha\wedge u)=\omega$, and put $u_*=J^*u$, where
$J(x,v)=(x,-v)$. Set $f=\Delta|\omega|^2$ with CDDP's nonpositive
Laplacian, and choose $b=Q_4f$. Their Theorem 5 and (5.6) give
\[
 S:=\int_M(\pi^*b)\alpha\wedge du\wedge du_*
      =\frac16\|f\|_2^2>0.
\]
Strict positivity follows from Proposition A.1. The pairing formulas
(3.55)--(3.57), (3.67), and their fiber orientation give
$\int_M\alpha\wedge d\alpha\wedge u\wedge u_*=\|\omega\|_2^2/(4\pi)$.
Equations (4.22) and (4.38) annihilate both the closed row and column of
the compressed first variation and give the remaining entry $iS$.
Since CDDP uses $P=-iL=iB$, the variation of $B$ is
$\operatorname{diag}(0,r)$, where
\[
 r=\frac{2\pi}{3\|\omega\|_2^2}\|f\|_2^2>0.
\]
This is CDDP Theorem 3's nondegeneracy condition when $b_1=1$;
that theorem and Corollary 4.1 give the asserted zero multiplicities.
Small conformal perturbations remain negatively curved.
\end{proof}

\section{The finite resonance divisor}
Write $c=B|_{U^0}$, $A_1=B|_{U^1}$, and $A_2=B|_{U^2}$.
All spaces below refer to the fixed cut $\gamma$.

\begin{lemma}\label{coefficients}
On one fixed complex disk $|z|<\rho$, the finite characteristic
coefficients are holomorphic in $z$ and continuous in $\tau$, locally
uniformly in $z$. They are even in $z$ and real on the real locus.
There are continuous, character-holomorphic functions $C,T,A$ such that
\begin{align}
 c=z^2C,&\qquad t:=\tr A_1-c=t(0,\tau)+z^2T,
      &C(0,0)=T(0,0)=-\kappa/2,\label{factors}\\
 \det(s-A_1)&=(s-c)(s-t),&
 \det(s-A_2)=(s-c)(s-t)(s^2-e_1s+e_2),\label{polys}\\
 e_1=z^2A,&\qquad A(0,0)=0,&
 e_1(z,0)=0,\quad e_2(z,0)=-c_0(z)^2.\label{middle}
\end{align}
Also $e_1(0,\tau)=e_2(0,\tau)=0$.
\end{lemma}
\begin{proof}
Chaubet--Dang Appendix B gives common anisotropic spaces
$\mathcal H_j$ on a finite list of levels, uniform same-level contour
resolvent bounds, and metric continuity $\mathcal H_{j+1}\to\mathcal H_j$.
The connection perturbation is bounded multiplication:
$L_{z,\tau}=L_{0,\tau}+zV_\tau$, $V_\tau=i\pi^*\omega(X_\tau)$.
Consequently, on a single disk independent of $\tau$,
\begin{equation}\label{neumann}
 R_{z,\tau}=R_{0,\tau}(\Id+zV_\tau R_{0,\tau})^{-1}
\end{equation}
is a uniformly convergent Neumann series. Bounded perturbation preserves
the graph domain at each fixed metric; no common graph domain across
metrics is required. The resolvent identity gives locally uniform metric
continuity with one lost level. Contour integration gives these properties
for $\Pi$ and $B\Pi$. Fixed smooth sources $f_j$ and tests $\ell_i$ with
invertible pairing matrix $(\ell_i\Pi f_j)$ produce finite frames and
actual compressed matrices. Insert $\Id-\alpha\wedge\iota_X$ for
horizontal degrees. Nearby projectors have the ranks \eqref{ranks}.

Pairing every orbit with its reversal proves evenness of each horizontal
Euler determinant. The two holonomies are $e^{\pm iz\int_\gamma\omega}$,
and a symplectic return map and its inverse have the same exterior traces
and $|\det(\Id-P_\gamma)|$. For real $z$ the paired coefficients are
real. Continuation proves the corresponding identities for the cut
characteristic polynomials.

Put $D=d_z-\alpha\wedge L$ and $W=d\alpha\wedge$. Contact calculus gives
\begin{equation}\label{contact}
 [B,D]=0,\qquad D^2=WB.
\end{equation}
For real nonzero $\theta$, scalar zero exclusion
\cite[Lemma 7.4]{DGRS} gives $c\ne0$. Applying \eqref{contact} to its
scalar eigenstate shows that $c$ occurs in $A_1$, proving the first
factorization in \eqref{polys}; it extends by holomorphy. As $c(0,\tau)=0$,
evenness and Cauchy's formula give \eqref{factors}. Explicitly, if
$h(0,\tau)=0$ is even, then
\[
 \frac{h(z,\tau)}{z^2}=\frac1{2\pi i}\int_{|w|=a}
             \frac{h(w,\tau)}{w^2(w-z)}\,dw,\qquad |z|<a<\rho.
\]
This proves continuity of the quotient, including at $z=0$.
Equation \eqref{c0} supplies its origin values, so $c\ne0$ for all
small complex $z\ne0$.

When $A_1$ is invertible, $D:U^1\to U^2$ is injective by
\eqref{contact} and the Lefschetz isomorphism $W:U^1\to U^3$.
Its characteristic polynomial therefore divides that of $A_2$.
For each fixed $\tau>0$, $t(0,\tau)>0$; hence the invertible locus
is dense in $z$. At $\tau=0$ use \eqref{hyper}. Monic division extends
the second factorization in \eqref{polys} everywhere. Lemma~\ref{axes}
gives \eqref{middle} and the zero-twist identities. Applying the same
Cauchy formula to $e_1$ yields $e_1=z^2A$ with $A(0,0)=0$.
\end{proof}

The horizontal alternating determinant and Lefschetz symmetry give
\begin{equation}\label{divisor}
 Z(s,z,\tau)=G(s,z,\tau)
          \frac{(s-c)(s^2-e_1s+e_2)}{s-t},
\end{equation}
where the cut remainder $G$ is holomorphic and nonzero in $s$ inside
$\gamma$; see \cite[(6.4)--(6.5)]{CD}.

\begin{lemma}\label{normalization}
For real nonzero small $\theta$,
\begin{equation}\label{Gvariation}
 \frac{G(0,\theta,\tau)}{G(0,\theta,0)}
   =\exp\left(\int_0^\tau H(\theta,v)\,dv\right),\qquad
 H=\str_{C^\bullet}\bigl(\Pi\alpha\wedge\iota_{\dot X}\bigr).
\end{equation}
The finite trace $H$ is uniformly bounded through $\theta=0$.
\end{lemma}
\begin{proof}
Pullback acyclicity follows from the sphere-bundle spectral sequence.
The full zero complex computes twisted de Rham cohomology, and every
nonzero spectral block has contraction $\iota_XL^{-1}$. Thus the full
cut complex is acyclic for fixed $\theta\ne0$.

The cut projector commutes with $d_\theta$, $\iota_X$, $\alpha\wedge$,
and $W$, so contact chirality $\Gamma$ preserves this arbitrary cut.
Transport by $\Pi_\tau|_{C_{\theta,\tau_0}}$ fixes the differential.
The identities $\Pi\dot\Pi\Pi=0$ and $[\Gamma,\Pi]=0$ cancel projector
terms in the chirality-torsion derivative. In dimension five, direct
differentiation of \cite[Definition 6.1]{CD} gives
\[
 \partial_\tau\log\tau(C^\bullet,\Gamma)
       =2h_0-2h_1+h_2=-H,\qquad
 h_k=\tr_{U^k}\bigl(\Pi\alpha(\dot X)\bigr).
\]
Here the $d\dot\alpha$ terms cancel between consecutive degrees;
Lefschetz gives $h_3=h_1$, $h_4=h_0$, and
$H=-h_0+h_1-h_2+h_3-h_4$. This is the calculation of
\cite[Lemma 7.3]{CD}; it uses no restriction of the cut to zero
resonances. Formula (6.6) and contact deformation invariance give
\[
 \tau_{\rm dyn}=(-1)^{Q_{\rm cut}}G(0,\theta,\tau)
                    \tau(C^\bullet,\Gamma),\qquad Q_{\rm cut}=3,
\]
which proves \eqref{Gvariation}.

For uniformity use the fixed sources and tests above. If
$K_\tau=\alpha_\tau\wedge\iota_{\dot X_\tau}$, its compressed matrix
is $F^{-1}J$, where
\[
 F_{ij}=\ell_i(\Pi f_j),\qquad
 J_{ij}=\ell_i(\Pi K_\tau\Pi f_j).
\]
Take sources in $\mathcal H_2$, use continuity
$\mathcal H_2\to\mathcal H_1$ for the inner projector and
$\mathcal H_1\to\mathcal H_0$ for the outer one, and boundedness of the
smooth order-zero map $K_\tau$. Then $F^{-1}J$ is continuous through
$\theta=0$. In particular the ratio in \eqref{Gvariation} tends to one
uniformly as $\tau\to0$. No scalar torsion is evaluated at the
nonacyclic character.
\end{proof}

\section{The crossing and the counterexample}
\begin{proof}[Proof of Theorem~\ref{main}]
For real $\theta\ne0$ small, $t(\theta,0)=c_0(\theta)<0$, whereas
$t(0,\tau)=r\tau+o(\tau)$. By \eqref{factors}, $|T|\le K$ and
$t(0,\tau)\ge r\tau/2$ near zero. At $\tau_1=4K\theta^2/r$ the
companion root is positive. Continuity therefore supplies a crossing
\begin{equation}\label{cross}
 t(\theta,\sigma_\theta)=0,\qquad
 0<\sigma_\theta\le4K\theta^2/r,\qquad
 \sigma_\theta/\theta^2\longrightarrow\kappa/(2r).
\end{equation}
At that crossing $c\ne0$ and the horizontal one-zero eigenvalue is simple.

Let $V^k$ be the horizontal generalized zero spaces. Their dimensions
are $(0,1,m,1,0)$. The full zero complex is acyclic. In the contact
splitting its differential is
\[
 d_\theta(u+\alpha\wedge v)
       =Du+Wv+\alpha\wedge(Lu-Dv).
\]
Exactness in degrees one and four makes $D:V^1\to V^2$ injective
and $D:V^2\to V^3$ surjective. Moreover $D^2=0$ on $V^1$, so
$m\ge2$. Since the rank-four middle cut contains the nonzero root $c$,
$m\le3$. Thus $e_2=0$ at the crossing, and
$Q(s)=s(s-e_1)$, with $e_1=0$ precisely when $m=3$.

Specialize the spectral fiber in \eqref{divisor} and cancel its common
$s$ factor. This gives the holomorphic function
\begin{equation}\label{crossvalue}
 Z(s,\theta,\sigma_\theta)
   =G(s,\theta,\sigma_\theta)(s-c)(s-e_1),\qquad
 Z(0,\theta,\sigma_\theta)=G(0,\theta,\sigma_\theta)c e_1.
\end{equation}
For $m=3$ this has a simple zero; for $m=2$ its value at zero is nonzero.
On the hyperbolic axis, by \eqref{hyper} and \eqref{middle},
\[
 Z(0,\theta,0)=-G(0,\theta,0)c_0(\theta)^2\ne0.
\]
Consequently, in both cases,
\begin{equation}\label{exactratio}
 \frac{Z(0,\theta,\sigma_\theta)}{Z(0,\theta,0)}
 =-\frac{G(0,\theta,\sigma_\theta)}{G(0,\theta,0)}
     \frac{c(\theta,\sigma_\theta)}{c_0(\theta)}
     \frac{e_1(\theta,\sigma_\theta)}{c_0(\theta)}
 \longrightarrow0.
\end{equation}
The three factors tend respectively to $1,1,0$: use
Lemma~\ref{normalization}, then \eqref{factors}, then
$e_1=\theta^2A$ with $A(0,0)=0$, together with \eqref{c0}.

Fried's hyperbolic theorem \cite{Fried} gives
$|Z(0,\theta,0)|=T_{\chi_\theta}(\Sigma)^2$.
Acyclic unitary analytic torsion is metric-independent on the closed
odd-dimensional base. Choose any nonzero $\theta_j\to0$ and put
$g_j=g_{\sigma_{\theta_j}}$, $\chi_j=\chi_{\theta_j}$.
Equation \eqref{exactratio} proves \eqref{ratio}; \eqref{cross} gives
smooth convergence of the negatively curved metrics.
\end{proof}

\paragraph{Specialization.}
The value in \eqref{crossvalue} is the ordinary value of the meromorphic
function for the crossing metric. Continuing nonresonant values instead
would amount to an additional cancellation in the parameter as well as
in $s$. Meromorphic continuation does not supply that cancellation.
Chaubet--Dang's dynamical torsion includes the refined torsion of the
zero complex \cite[(6.7)]{CD}; its invariance is consistent with the
change of the ordinary specialized value.

\begin{thebibliography}{9}
\bibitem{BO} U. Bunke and M. Olbrich,
\emph{Theta and zeta functions for odd-dimensional locally symmetric spaces
of rank one}, \href{https://arxiv.org/abs/dg-ga/9407012}{arXiv:dg-ga/9407012},
1994, Proposition 4.7.
\bibitem{CDDP} M. Ceki\'c, B. Delarue, S. Dyatlov and G. P. Paternain,
\emph{The Ruelle zeta function at zero for nearly hyperbolic 3-manifolds},
Invent. Math. \textbf{229} (2022), 303--394,
\href{https://arxiv.org/abs/2009.08558}{arXiv:2009.08558v3}.
\bibitem{CD} Y. Chaubet and N. V. Dang,
\emph{Dynamical torsion for contact Anosov flows},
Anal. PDE \textbf{17} (2024), 2619--2681,
\href{https://arxiv.org/abs/1911.09931}{arXiv:1911.09931v2}.
\bibitem{DGRS} N. V. Dang, C. Guillarmou, G. Rivi\`ere and S. Shen,
\emph{The Fried conjecture in small dimensions},
Invent. Math. \textbf{220} (2020), 525--579,
\href{https://arxiv.org/abs/1807.01189}{arXiv:1807.01189}.
\bibitem{Fried} D. Fried,
\emph{Analytic torsion and closed geodesics on hyperbolic manifolds},
Invent. Math. \textbf{84} (1986), 523--540.
\bibitem{MS} D. Margalit and S. Spallone,
\emph{A homological recipe for pseudo-Anosovs},
\href{https://margalit.droppages.net/papers/homcrit.pdf}{author manuscript},
Propositions 1--2.
\bibitem{Thurston} W. P. Thurston,
\emph{Hyperbolic structures on 3-manifolds, II: surface groups and
3-manifolds which fiber over the circle},
\href{https://arxiv.org/abs/math/9801045}{arXiv:math/9801045}.
\bibitem{Zw} M. Zworski,
\emph{Mathematical study of scattering resonances},
Bull. Math. Sci. \textbf{7} (2017), 1--85,
\href{https://doi.org/10.1007/s13373-017-0099-4}{doi:10.1007/s13373-017-0099-4},
Conjecture 9.
\bibitem{Crossing} blawton,
\emph{Resonance-cluster analysis for a proposed counterexample to Fried's
conjecture}, conditional Lean formalization, 2026,
\href{https://github.com/blawton/fried-conj-lean/tree/610d659ea15af869553f6e3e79c432880d9b0ab0}
{\texttt{fried-conj-lean}}, revision \texttt{610d659};
\href{https://palomar-registry.org/entry?id=PALOMAR-2026-09-17-000006\&version=1}
{Palomar entry PALOMAR-2026-09-17-000006}.
\end{thebibliography}
\end{document}
